% !TeX root = ../教材题库_全部条目_转写工作稿.tex
% Source: XJ-B1-C05-PR-P158-Y0346
% 第5章 / 5.1 任意角和弧度制 / 教材 p.158 / PDF 第165页
% OCR未校对：公式与文字均可能有误；未完成数学TeX转写；禁止直接发给学生
% 原题及相邻原解截图（核对用）：../题目原图/XJ-B1-C05-PR-P158-Y0346-01.webp
% 本文件只在转写完成后作为题干使用；OCR 条目尚不是可靠题干。
\begin{bankitem}{p.158 练习 第3题}{第5章 / 5.1 任意角和弧度制 / 教材 p.158 / PDF 第165页}{XJ-B1-C05-PR-P158-Y0346}{OCR未校对：公式与文字均可能有误；未完成数学TeX转写；禁止直接发给学生}
\begin{ocrdraft}
[教材p.158 / PDF第165页]
\par
3．在0。、360。内找出与下列各角终边相同的角，并指出它们是第几象限角：
\par
一545。30'；
\par
（1)3290。；
\par
（2）
\par
（3）
\par
一210。；
\par
（4）
\par
一1300。．
\end{ocrdraft}
\end{bankitem}
